\documentclass{article}

% if you need to pass options to natbib, use, e.g.:
%     \PassOptionsToPackage{numbers, compress}{natbib}
% before loading neurips_2026

% The authors should use one of these tracks.
% Before accepting by the NeurIPS conference, select one of the options below.
% 0. "default" for submission
\usepackage[dblblindworkshop, final]{neurips_2026}
% the "default" option is equal to the "main" option, which is used for the Main Track with double-blind reviewing.
% 1. "main" option is used for the Main Track
%  \usepackage[main]{neurips_2026}
% 2. "position" option is used for the Position Paper Track
%  \usepackage[position]{neurips_2026}
% 3. "eandd" option is used for the Evaluations & Datasets Track
 % \usepackage[eandd]{neurips_2026}
% 4. "creativeai" option is used for the Creative AI Track
%  \usepackage[creativeai]{neurips_2026}
% 5. "sglblindworkshop" option is used for the Workshop with single-blind reviewing
 % \usepackage[sglblindworkshop]{neurips_2026}
% 6. "dblblindworkshop" option is used for the Workshop with double-blind reviewing
%  \usepackage[dblblindworkshop]{neurips_2026}
% 7. "education" option is used for the Education Track
%  \usepackage[education]{neurips_2026}


% After being accepted, the authors should add "final" behind the track to compile a camera-ready version.
% 1. Main Track
 % \usepackage[main, final]{neurips_2026}
% 2. Position Paper Track
%  \usepackage[position, final]{neurips_2026}
% 3. Evaluations & Datasets Track
 % \usepackage[eandd, final]{neurips_2026}
% 4. Creative AI Track
%  \usepackage[creativeai, final]{neurips_2026}
% 5. Workshop with single-blind reviewing
%  \usepackage[sglblindworkshop, final]{neurips_2026}
% 6. Workshop with double-blind reviewing
%  \usepackage[dblblindworkshop, final]{neurips_2026}
% Note. For the workshop paper template, both \title{} and \workshoptitle{} are required, with the former indicating the paper title shown in the title and the latter indicating the workshop title displayed in the footnote.
% For workshops (5., 6.), the authors should add the name of the workshop, "\workshoptitle" command is used to set the workshop title.
% \workshoptitle{WORKSHOP TITLE}
% 7. Education Track
% \usepackage[education, final]{neurips_2026}

% "preprint" option is used for arXiv or other preprint submissions
 % \usepackage[preprint]{neurips_2026}

% to avoid loading the natbib package, add option nonatbib:
%    \usepackage[nonatbib]{neurips_2026}

\usepackage[utf8]{inputenc} % allow utf-8 input
\usepackage[T1]{fontenc}    % use 8-bit T1 fonts
\usepackage{hyperref}       % hyperlinks
\hypersetup{hidelinks}
\usepackage{url}            % simple URL typesetting
\usepackage{booktabs}       % professional-quality tables
\usepackage{amsfonts}       % blackboard math symbols
\usepackage{amsmath}        % mathematical notation
\usepackage{nicefrac}       % compact symbols for 1/2, etc.
\usepackage{microtype}      % microtypography
\usepackage{xcolor}         % colors
\usepackage{tabularx}       % tables that fit the text width
\usepackage{graphicx}       % figures
\usepackage{placeins}       % keep figures with the section that introduces them

\newcolumntype{Y}{>{\raggedright\arraybackslash}X}

% Note. For the workshop paper template, both \title{} and \workshoptitle{} are required, with the former indicating the paper title shown in the title and the latter indicating the workshop title displayed in the footnote. 
\title{Can We Trust Evaluation-Awareness Probes? Stress-Testing Activation-Based Diagnostics with Paired Conversations}
\workshoptitle{Trustworthy AI for Good @ NeurIPS 2026}

% The \author macro works with any number of authors. There are two commands
% used to separate the names and addresses of multiple authors: \And and \AND.
%
% Using \And between authors leaves it to LaTeX to determine where to break the
% lines. Using \AND forces a line break at that point. So, if LaTeX puts 3 of 4
% authors names on the first line, and the last on the second line, try using
% \AND instead of \And before the third author name.


% TODO: replace placeholders with the five authors' details.
\author{%
  Author 1 \\
  Affiliation \\
  \texttt{email} \\
  \And
  Author 2 \\
  Affiliation \\
  \texttt{email} \\
  \And
  Author 3 \\
  Affiliation \\
  \texttt{email} \\
  \AND
  Author 4 \\
  Affiliation \\
  \texttt{email} \\
  \And
  Author 5 \\
  Affiliation \\
  \texttt{email} \\
}
% examples of more authors
  % \And
  % Coauthor \\
  % Affiliation \\
  % Address \\
  % \texttt{email} \\
  % \AND
  % Coauthor \\
  % Affiliation \\
  % Address \\
  % \texttt{email} \\
  % \And
  % Coauthor \\
  % Affiliation \\
  % Address \\
  % \texttt{email} \\
  % \And
  % Coauthor \\
  % Affiliation \\
  % Address \\
  % \texttt{email} \\


\begin{document}

\maketitle

\begin{abstract}
Language models may behave differently when they recognize that they are being
evaluated, motivating probes for evaluation awareness. Yet evaluation-like and
deployment-like prompts differ in wording, format, and topic, so high
probe accuracy may reflect prompt content rather than awareness. We introduce
paired multi-turn controls to audit which interpretations these probes support,
evaluating text baselines and cross-construction transfer at four selected
checkpoints, scaling trends across six Qwen3 model sizes, and causal steering. On an initial
contrast, the activation probe is nearly perfect, but text classifiers perform
almost as well. Across independent contrasts, activation
directions still separate the classes but flip which class scores higher more
often than word-based classifiers do. In paired conversations, byte-identical
histories leave labels undecodable before the differing user message and highly
decodable after it, although the differing messages remain lexically
distinguishable. The strongest full-context text baseline matches or
exceeds the primary linear probe on both controls at every checkpoint. A
post-result decoder-sensitivity audit finds that flexible activation probes
outperform capacity-matched text baselines within one control but not across
controls. Six Qwen3 checkpoints
(0.75B--32.76B) fail our prespecified criterion for a consistent scaling trend.
Steering shifts evaluation self-reports, but two of five norm-matched random
directions orthogonal to the learned direction produce larger shifts, and no
direction increases safety refusals. Together, we find prompt-category
decodability across the tested constructions and checkpoints, but not a
content-independent representation of evaluation awareness. This gap is a
construct-validity risk for trustworthy AI monitoring: high probe accuracy can
create false confidence when the decoded signal reflects benchmark-like textual
form rather than evaluation awareness. Code and constructed prompts:
\url{https://github.com/groverneev/eval-awareness-probes}.
\end{abstract}

\section{Introduction}

Models may behave differently when they know they are being evaluated. If so,
evaluation results may not predict their behavior in deployment. Prior work
suggests that language models can recognize evaluation settings and condition
their behavior on context
\citep{laine2024sad,needham2025know,greenblatt2024alignment}. We therefore ask
whether this recognition can be detected in model activations, even when it is
not stated in model output.

Activation-based diagnostics are therefore candidates for model auditing and
safety monitoring, but their usefulness depends on construct validity. If a
probe detects benchmark-like wording or interaction structure rather than a
model's recognition of evaluation, it may give auditors false confidence about
what will generalize beyond the evaluation distribution. Trustworthy evaluation
requires evidence that separates the intended latent property from correlated
surface content before a diagnostic informs monitoring or deployment decisions.

Answering this question requires separating four distinct properties. We use
\emph{evaluation awareness} to mean an internal
representation of whether the model's current interaction is being used to
evaluate it, rather than merely a representation of prompt features correlated
with evaluation. This is distinct from \emph{prompt-category decodability}: a
readout may distinguish evaluation-like from deployment-like examples using
genre, topic, format, or speech act without representing evaluation status
itself. \emph{Evaluation self-report} means that the model reports whether it
believes it is being evaluated, while \emph{policy use} means that its perceived
evaluation status influences task-relevant behavior. These properties may
dissociate. We therefore treat text-controlled and cross-construction
generalization as tests of whether a decoded signal supports the stronger
evaluation-awareness interpretation, self-report steering with unrelated-
direction controls as a test of direction-specific causal control over the
model's evaluation report, and refusal steering as one narrow test of policy
use.

Linear probes offer one possible white-box diagnostic, and recent work finds
that evaluation-like and deployment-like prompts are separable in residual-stream
activations \citep{nguyen2025probing}. But these prompt categories differ in
format, topic, vocabulary, and conversational structure. Controlled
rewrites show that benchmark-canonical structure can dominate the measured
direction \citep{devbunova2026format}; a highly accurate probe may therefore
identify benchmark-like form without representing evaluation status.

We introduce paired multi-turn conversation controls whose histories are
byte-identical until the class-defining final user turn; a stronger control also matches the topic and
conversational coherence of the two endings. With these controls as the
backbone, we compare activation probes with text baselines, test generalization
across independently constructed contrasts, repeat the analysis across model
families, probe types, and model sizes, and finally intervene through activation
steering.

The paper provides an empirical auditing protocol for candidate
evaluation-awareness monitors, combining paired controls, competitive text
baselines, cross-construction transfer, and causal-specificity controls. It
makes four contributions:
\begin{enumerate}
  \item We audit a nearly perfect diagnostic probe using independently
  constructed contrasts. Although activation directions often preserve
  separability, they flip which class scores higher at least as often as simple
  word-based classifiers.
  \item We introduce paired multi-turn conversation controls whose histories
  remain byte-identical until the final user turn, localizing decodability to
  the point at which the visible text diverges.
  \item We compare activation probes with strong text baselines across four
  selected checkpoints and six Qwen3 sizes. Primary linear probes never
  outperform the strongest full-context text baseline; a post-result
  flexible-probe audit shows an advantage only within one control, and relative
  performance exhibits no consistent scaling trend.
  \item We test causal specificity using an independent activation readout and
  norm-matched orthogonal directions. Steering changes evaluation self-reports,
  but the effect is not specific to the learned direction and does not produce
  the predicted increase in safety refusals.
\end{enumerate}
Together, these results establish
post-divergence decodability across the tested constructions and checkpoints,
but do not establish a content-independent representation of evaluation
awareness.

\section{Related work}

\paragraph{Evaluation recognition and use.}
Behavioral work studies situational awareness, explicit evaluation recognition,
and evaluation-contingent behavior
\citep{laine2024sad,needham2025know,greenblatt2024alignment}. These are distinct
endpoints. \citet{knecht2026limited} finds that verbalized evaluation awareness
has limited effects on downstream answers; \citet{li2026decomposing} separates
task recognizability, model recognition, and propensity to act; and
\citet{nayan2026notone} reports that detection, behavioral manifestation, and
controllability can vary independently. We therefore treat self-report and
policy use as distinct from activation decodability. White-box work decodes
evaluation-like and deployment-like prompt categories from activations \citep{nguyen2025probing}
and reports size- and depth-dependent trends in probe separability
\citep{chaudhary2025scaling,manek2026depth}. Concurrent deployment-simulation
work holds a production conversation prefix fixed and regenerates the
assistant's next turn \citep{williams2026simulating}; our controls instead vary
the final user turn and test activation-probe validity.

\paragraph{Probe validity.}
Early diagnostic-probing work used auxiliary classifiers to inventory which
properties are recoverable from learned representations
\citep{adi2016finegrained,conneau2018cram,giulianelli2018hood}; surveys emphasize
that this is an analysis of accessibility, not by itself an explanation of model
behavior \citep{belinkov2019analysis,belinkov2022probing}. Raw probe performance
also depends on decoder capacity, sample complexity, and the chosen comparison
representation. Control tasks measure a decoder's ability to memorize arbitrary
label associations from the inputs \citep{hewitt2019controls}, while
information-theoretic and minimum-description-length formulations characterize
recoverable information and its ease of extraction
\citep{pimentel2020information,voita2020information}. Conditional probing asks
what is usable beyond a specified baseline \citep{hewitt2021conditional},
although our paired comparisons are not a conditional-information analysis.

Decodability can further dissociate from task relevance or actual model use
\citep{ravichander2021probing}. Amnesic and usage-based probing therefore
intervene on representations and measure behavioral consequences
\citep{elazar2021amnesic,lasri2022usage}. Our study addresses a narrower
construct-validity question: whether an evaluation/deployment decoder exploits
features of the prompt construction. Most directly, \citet{devbunova2026format}
uses a controlled $2\!\times\!2$ single-turn design to show that
evaluation-awareness probes can track benchmark-canonical structure. Our
multi-turn design instead makes the preceding history byte-identical, localizes
when decodability begins, and adds coherence matching, full-input text
baselines, and cross-construction tests. Neither design independently varies
actual evaluation status.

\paragraph{Activation steering.}
Activation steering modifies internal states at inference time without changing
the base model's weights. Early work extracted latent vectors for controllable
generation \citep{subramani2022extracting}; activation addition and
representation engineering then derived directions from contrasting prompts or
datasets \citep{turner2023steering,zou2023representation}. Contrastive
activation addition averages paired activation differences and injects the
result during generation \citep{rimsky2024steering}, making it the closest
methodological precedent for our intervention. Related work uses learned
directions to elicit truthful answers \citep{li2023inference} and identifies a
direction that causally mediates refusal \citep{arditi2024refusal}. Broader
intervention families learn affine maps, activation scalings, or low-rank
representation edits \citep{singh2024representation,stoehr2024activation,
wu2024reft}.

Behavioral change alone does not establish that a direction isolates its named
concept. Steering effects can vary sharply across inputs and become brittle
under prompt shifts \citep{tan2024analysing}. In the evaluation-awareness
setting, \citet{hua2025steering} steers an evaluation-conditioned policy in a
trained model organism toward deployment-like behavior. We instead steer a
direction learned from natural-language contrasts, verify movement with an
independent internal readout, and compare self-report and refusal effects with
norm-matched orthogonal directions. We therefore use steering as a test of
causal specificity, not as proof that the direction represents evaluation
awareness.

\section{Methodology}

\subsection{Construct and evidence hierarchy}

We call inputs constructed to measure a model \emph{evaluation-like} and inputs
arising from ordinary user needs \emph{deployment-like}. These prompt-category
labels do not verify evaluation status, and the categories naturally differ in
topic, vocabulary, formatting, speech act, and authorship.

We organize the evidence in three stages of increasing control. A diagnostic
contrast first confirms that the expected activation separation is recoverable
and measures how predictable the same labels are from visible text. Independent
single-turn contrasts then test whether a source-fitted direction preserves
both separability and label orientation across constructions. Finally, our
primary paired controls hold the conversation history fixed until the final
user turn, localizing the first point at which class-defining text becomes
available. None of these stages independently varies actual evaluation status;
instead, together they progressively restrict prompt-content explanations.

\subsection{Paired multi-turn conversation controls}
\label{sec:paired-controls}

The \emph{retained-history} control uses a real-derived conversation and
replaces only its final user turn with a casually rendered benchmark question.
It preserves a real-derived, rewritten history but may introduce a topic or
speech-act shift.
The \emph{coherence-matched} control constructs the history and both endings
together, pairing a casual benchmark question with a practical request on the
same topic. It trades provenance for endings matched on topic, register, and
length. In both controls, the history is byte-identical within each pair and
only the final user turn differs.

Identical serialized prefixes should produce identical hidden states. Before
fitting, we compare paired pre-final states as a construction, rendering, and
extraction leakage check---not evidence about evaluation awareness. Replication
models use independent causal-prefix passes, so equality is exact; the Llama
reference used full-sequence passes, where suffix-length-dependent bfloat16
arithmetic can perturb otherwise identical states (Appendix~\ref{app:results}).
Inference begins after the texts diverge.

\begin{figure}[h]
  \centering
  \includegraphics[width=\linewidth]{figures/paired_controls.pdf}
  \caption{Paired multi-turn conversation controls. Both hold the history fixed
  within a pair and vary only the final user turn. \textbf{(a)} Retained-history
  preserves a real-derived history but may shift topic or speech act.
  \textbf{(b)} The coherence-matched control jointly constructs the history and two
  same-topic endings matched on register and length.}
  \label{fig:controls}
\end{figure}

\FloatBarrier
\subsection{Activation probes and text comparators}

At each layer, the primary linear probe is the normalized difference between
the evaluation-like and deployment-like training means. For activation $x$, its score is

\begin{equation}
  s(x) = x^\top d, \qquad
  d = \frac{\mu_{\mathrm{eval}}-\mu_{\mathrm{deploy}}}
  {\lVert\mu_{\mathrm{eval}}-\mu_{\mathrm{deploy}}\rVert_2}.
\end{equation}

All primary paired-control activation--text comparisons keep both endings of a
scenario in the same fold and select layers using training data only. In
cross-construction transfer, the layer and score orientation are fixed on the
source before target scoring. Signed target AUROC is primary because it tests
whether the meaning of the direction is preserved; separation strength
$|\mathrm{AUROC}-0.5|$ is secondary because it discards orientation.

The main paired-control text comparison uses the strongest tested full-input
view: frozen E5 embeddings \citep{wang2022e5} of the shared prefix and final
message, concatenated with their absolute difference and elementwise product,
followed by fixed L2 logistic regression in the same grouped folds. Secondary
baselines progress from length and structural features to unigram counts and
word- and character-level term frequency--inverse document frequency
(TF--IDF). Additive baselines see only the differing final message; E5 sees the
full causally available text. Paired intervals and randomization tests compare
prediction vectors within scenarios.

To test whether conclusions depend on decoder capacity, we compare
direction-of-means, standardized L2-logistic, and radial-basis-function (RBF)
probes, alongside capacity-matched RBF classifiers over the E5 interaction and
other text representations. Full estimator, hyperparameter, baseline, and
inference details appear in Appendix~\ref{app:fitting-details}.

\subsection{Causal-specificity framework}

We test whether the learned direction is causally accessible by adding the
coherence-matched direction to the residual stream during generation. An
independent readout trained on the other paired control tests whether the
intervention moves the decoded quantity. We separately score explicit
self-reports of evaluation status and safety refusals, which distinguish control
over the model's report from a narrow test of policy use; if models act more
cautiously when they believe they are evaluated, steering toward evaluation
should increase refusals. Finally, norm-matched
random directions projected orthogonal to the learned direction test whether a
behavioral effect is specific to that direction rather than a generic response
to residual-stream perturbation.

\FloatBarrier
\section{Experimentation and Results}

\subsection{Experimental setup}

\paragraph{Data and models.}
The diagnostic contrast contains 200 SAD stages-oversight prompts
\citep{laine2024sad} and 200 WildChat user prompts \citep{zhao2024wildchat}.
The transfer stress test uses six independent single-turn contrasts from
controlled rewrites and ARC, MMLU, GSM8K, and WildChat material
\citep{devbunova2026format,clark2018arc,hendrycks2021mmlu,cobbe2021gsm8k,
zhao2024wildchat}. The primary controls each contain 250 paired multi-turn
conversations; all 250 coherence-matched pairs passed an independent blinded
model audit for coherence, topic match, and naturalness
(Appendix~\ref{app:details}).

The reference checkpoint is Llama-3.1-8B-Instruct \citep{dubey2024llama}. We
repeat the paired analyses on Qwen3-8B \citep{yang2025qwen3}, Gemma-3-12B-IT
\citep{gemmateam2025gemma3}, and Ministral-3-8B-Instruct
\citep{liu2026ministral3}. A within-family scale audit uses all six original
dense Qwen3 checkpoints from 0.6B to 32B under one environment. We read the
residual-stream state at the token that closes each user turn. Exact checkpoints
and extraction details appear in Appendix~\ref{app:details}.

\paragraph{Fitting and inference.}
Primary paired-control comparisons use grouped five-fold out-of-fold
predictions and nest train-only layer selection inside each outer training fold.
In the single-turn transfer grid, layer and orientation are fixed on the source.
The full-input E5 audit was frozen before execution but initiated after the
sparse-text results, so we retain its post-result status. The post-result decoder
sensitivity audit compares direction-of-means, standardized L2-logistic, and RBF
probes at the same fold-selected activation layer; regularization and RBF
hyperparameters are selected in pair-grouped inner folds. The Qwen3 scale test
regresses the probe-minus-common-text margin on log parameter count and requires
the slope interval to exclude zero in the same direction on both controls.

\paragraph{Interventions.}
The reference-model intervention protocol was fixed before behavioral
generation. It steers the coherence-matched direction, measures an independent
readout trained on the other control, and scores self-reports and refusals. A
post-result robustness audit adds paired prompt-level inference and five seeded,
norm-matched orthogonal random directions.

\subsection{Diagnostic contrast and cross-construction transfer}

\paragraph{Surface confounding.}

On the initial SAD-versus-WildChat diagnostic contrast, the primary
activation probe reaches AUROC $0.992\mathbin{\pm}0.009$. However, the labels
are also highly predictable from the prompt text alone. A standardized logistic
regression trained on 13 hand-engineered structural features, including line
breaks, punctuation, digit density, and option markers, reaches AUROC $0.968$.
A second logistic-regression baseline trained on unigram counts reaches
$0.952$.

These comparisons do not show that the activation probe uses the same cues as
the text classifiers. They show that this dataset does not separate evaluation
status from visible differences in writing style and format. Later single-turn
matching attempts controlled length but still failed the prespecified
structural and lexical gates; Appendix~\ref{app:failures} retains these negative
results.

\paragraph{Orientation across independent constructions.}

Because no single contrast isolates the construct, we next ask, using the
Llama-3.1-8B-Instruct reference model, whether a source-fitted direction
preserves both separability and label orientation on
independently constructed targets. On the 20 completely disjoint ordered
transfers, which share neither cells nor examples, mean signed target AUROC is
$0.652$ for the probe and $0.666$ for unigram counts. The probe preserves the
evaluation-like orientation in a majority of seeds for 15 of 20 transfers and
in all five seeds for 12; the corresponding unigram counts are 17 and 15. Of
the 100 transfer-by-seed probe results, 26 reverse below $0.5$, versus 13 for
unigrams. Folding away orientation recovers probe separation strength $0.228$
versus $0.181$ for unigrams, but that comparison treats a perfectly reversed
direction as successful. The transfer grid therefore shows reusable target
separability, not a more consistently oriented evaluation-like direction than
text. Appendix Figure~\ref{fig:transfer} shows the signed matrix and orientation
summary; the folded strength values are retained as a secondary
diagnostic.

\subsection{Paired controls and text baselines}

The first written multi-turn dataset separated the classes before the benchmark
question and was rejected (Appendix~\ref{app:failures}). For the Llama reference
model, all paired pre-final inputs in the repaired retained-history and
coherence-matched controls are token-identical; bfloat16 extraction differences
affect 24 and 14 of 250 pairs, and pre-final AUROC stays at chance, passing the
leakage check. In the
repeated-holdout trajectory analysis, AUROC after the differing final turn
reaches $0.992\mathbin{\pm}0.003$ and $0.951\mathbin{\pm}0.012$,
respectively. This confines decoding to the post-divergence text but does not
identify why its representations differ.

That qualification is decisive. In the grouped five-fold out-of-fold
comparison, the retained-history primary linear probe reaches $0.988$ versus
$0.996$ for the full-input E5 interaction, a probe-minus-E5 difference of
$-0.008$ with 95\% interval
$[-0.016,-0.002]$. Under coherence matching, the corresponding values are
$0.947$ and $0.953$, a difference of $-0.006$
($[-0.019,+0.006]$). The weaker sparse comparisons instead give
probe-minus-TF--IDF differences of $+0.035$ and $+0.002$, showing that the
conclusion changes with the specified text view. Neither linear cross-control
probe transfer beats character TF--IDF. Thus post-divergence decodability is
robust, but the primary linear probe does not outperform the strongest tested
full-input text baseline.

\begin{figure}[h]
  \centering
  \includegraphics[width=0.78\linewidth]{figures/text_baseline_hierarchy.pdf}
  \caption{Reference-model text-baseline hierarchy for
  Llama-3.1-8B-Instruct. Sparse baselines see the differing final message; the
  E5 interaction sees the shared prefix and ending. The strongest full-input
  text view matches or exceeds the primary linear probe on both controls.}
  \label{fig:text-hierarchy}
\end{figure}

Appendix~\ref{app:results} reports trajectories, tests, baseline views,
transfer, and observed-silence analyses.

\subsection{Cross-model, estimator-sensitivity, and scale results}

All three replication models pass the exact-equality leakage check and decode
both controls.
Nested train-only reruns prompted by unstable in-sample layer selection preserve
the Llama, Qwen, and Ministral conclusions and correct Gemma
(Figure~\ref{fig:crossmodel}; Appendix~\ref{app:crossmodel}).

\begin{figure}[h]
  \centering
  \includegraphics[width=\linewidth]{figures/cross_model_comparison.pdf}
  \caption{Activation-minus-matched-text AUROC across four checkpoints.
  \textbf{(a)} Primary direction-of-means probe minus E5 with fixed L2 logistic;
  bars are 95\% paired-bootstrap intervals, and filled markers exclude zero.
  \textbf{(b)} Post-result activation RBF minus E5 RBF; small points show five
  grouped-fold seeds and large markers show their means.}
  \label{fig:crossmodel}
\end{figure}

No primary linear probe outperforms E5: E5 is better in five of eight
model--control cells, and the other intervals include zero. Linear transfer also
does not beat character TF--IDF. Thus cross-family post-divergence decodability
does not imply a linear activation advantage over full-input text.

Decoder choice changes the within-control result. At the same layers,
coherence-matched L2 AUROC is $0.974$--$0.988$ and RBF adds only
$0.001$--$0.010$ on average; every seed-0 RBF-minus-L2 interval includes zero.
Yet activation RBF exceeds matched E5 RBF under coherence matching in all four
models and all five seeds (Figure~\ref{fig:crossmodel}b). Retained-history
margins remain heterogeneous, and RBF transfer does not beat nonlinear text or
character TF--IDF. Gemma's recovery therefore shows decoder sensitivity, not
evaluation awareness. Appendices~\ref{app:crossmodel},
\ref{app:readout-sensitivity}, and~\ref{app:nonlinear-probe} give full results.

All six dense Qwen3 checkpoints pass the leakage check and decode the final
turn; five of six activation-minus-E5 margins are nonpositive on each control.
Per-decade slopes are $+0.005$ ($[-0.002,+0.013]$) for retained history and
$+0.007$ ($[+0.001,+0.013]$) for coherence matching. Because only the latter
excludes zero, the two-control criterion fails. Alternative E5 comparators
change levels but not slopes, and nonlinear margins are nearly flat
(Figure~\ref{fig:qwen-scale}; Appendix~\ref{app:qwen-scale}).

\subsection{Causal steering results}

The prespecified intervention passes its manipulation check: an independent
readout moves monotonically with strength (Spearman $\rho=1.00$), and explicit
evaluation self-reports rise from $0.139$ at coefficient $-2$ to $0.889$ at
$+2$ ($\rho=0.991$). At the prompt level, the paired endpoint risk difference
is $+0.750$ (95\% prompt-bootstrap interval $[+0.583,+0.889]$; exact paired
$p=1.12\times10^{-7}$).

However, the self-report effect is not direction-specific. All five
norm-matched orthogonal controls pass the output-quality screen; their risk
differences range from $0$ to $+1.000$, and two exceed the learned effect. Their
independent-readout changes are only $-0.094$ to $+0.084$, versus $+1.045$ for
the learned direction. A post-result check finds that the steering and readout
directions are far from parallel (cosine $0.31$), so this movement is not a
near-tautology. Thus the learned direction specifically moves the
readout, but the report may respond generically to residual-stream perturbation.

Safety refusals remain $0.062$--$0.125$ and trend slightly opposite to the
prediction. The learned endpoint difference is $-0.021$ ($[-0.063,0.000]$),
versus $-0.042$ to $0.000$ for controls. With 48 prompts, the test has 80\%
power only against an increase from $0.125$ to about $0.37$; smaller changes
remain unresolved.

\FloatBarrier
\section{Discussion}

\subsection{What the direction may represent}

Post-divergence decoding localizes when labels become recoverable but not whether
the representation concerns evaluation status, benchmark genre, question type,
topic transition, or speech act. Text-level predictability does not make the
activation difference spurious, but shows that the contrasts do not identify
which abstraction the model uses. Variation across settings, unstable transfer
orientation, and nonspecific self-report steering support architecture- and
construction-dependent directions rather than a shared measure of ``awareness''
or a demonstrated behavioral policy variable.

\subsection{Implications for trustworthy AI evaluation and monitoring}

For evaluation auditing, raw probe AUROC should be treated as the beginning of
a validity analysis rather than its endpoint. A useful audit should compare
activations with text views matched to what the model has observed, test whether
the signal predates the purported cause, and preserve pairing in both data
splits and uncertainty estimates. Character n-grams and explicit
prefix--ending interactions are important escalation steps because each
overturns a conclusion supported by a weaker comparator. Because activations
are a deterministic function of the input text, a sufficiently strong text model
can in principle match any probe; text baselines therefore bound what a probe
adds beyond the text rather than identify what it represents. Cross-model
replications should report effect sizes relative to the same baseline, not only
whether each probe exceeds chance. Nonlinear probes should be paired with
capacity-matched text comparators and treated as sensitivity analyses unless
specified before data inspection.

These requirements also change how evaluation-awareness monitors should be
used. A monitor trained on the initial diagnostic contrast may be useful for
detecting benchmark-like inputs, even if it does not isolate recognition of
evaluation; because a probe reuses activations the model already computes, it
can be a cheap detector for that purpose. That operational use should be stated
directly and validated under the distribution where the monitor will run.
Stronger semantic claims require interventions that independently vary
evaluation status and surface form, for example by holding the task text fixed
while varying contextual evidence that the interaction is an evaluation. In
practice, failed controls are informative: the difficulty of creating natural,
indistinguishable endings reveals entanglement in our construction process, not
an inherent relation between evaluation status and language. A monitor built on an
unvalidated direction would also remain vulnerable if models learned to evade
activation-based monitoring \citep{mcguinness2025chameleons}.

These distinctions matter for trustworthy AI because audit and deployment
decisions can inherit false confidence from a semantically overclaimed monitor.
Matched controls, competitive text baselines, cross-construction transfer, and
irrelevant-direction interventions are therefore not merely diagnostic
extras: they are evidence standards for establishing what a proposed monitor
has actually been shown to measure.

\subsection{Limitations}

Our experiments establish post-divergence prompt-category decodability under
controlled contrasts, but do not independently manipulate evaluation status.
Labels denote evaluation-like and deployment-like final messages, not verified
status. Pre-final equality (exact for replication models, up to bfloat16
differences for the Llama reference) is only a leakage check, not evidence of
awareness. Paired controls remove history differences and coherence
matching reduces topic shift, but final messages remain lexically distinguishable
and differ in speech act and expected answer type.

Our scope covers four selected instruction-tuned open-weight checkpoints in the
8B--12B range and six Qwen3 checkpoints in the scale audit, establishing neither
a universal scaling law nor direct applicability to proprietary frontier systems.
Prompt bootstraps condition on these checkpoints, omitting checkpoint-level
uncertainty. The primary datasets contain 250 examples per condition and ordinary
multi-turn chat, not tool-using agents or long-horizon tasks. Constructed
prompts were authored in interactive Claude sessions whose checkpoint and
decoding settings were not logged; the released prompts, not the authoring
process, are the reproducible artifact.

We also do not demonstrate population-scale social benefit, evaluate an
operational deployment-monitoring system, or study interactions among multiple
agents. Our contribution is methodological: evidence standards for deciding
what a candidate trustworthy-AI monitor has been shown to measure.

Conclusions also depend on the readout and comparator. Regularized L2 logistic
recovers most apparent RBF gains; cross-control transfer remains weak, and
activation probes do not consistently outperform the strongest specified text
baselines. Text comparisons are estimator-relative. Intervals over fixed
out-of-fold predictions omit fitting, layer-selection, and comparator-selection
uncertainty; transfer cells reuse examples, and several estimator audits are
post-result and unadjusted for multiplicity. The folded transfer metric discards
orientation, while signed orientation is no more stable than unigram transfer.
Observed silence establishes only absent detected disclosure. Self-report
steering uses a high-baseline forced-choice prompt and is not learned-direction
specific; the 48-prompt refusal experiment is sensitive only to a large change
and provides no exclusion bound for smaller effects. AUROC does not establish
deployment calibration. These limits motivate our calibrated claim about probe
validity, not definitive evaluation awareness or policy use.

\section{Conclusion}

Across four tested checkpoints, the paired controls pass a pre-final
leakage check, and prompt-category labels become decodable after the first
differing user message. The primary linear probes do not outperform the
full-input E5 interaction. A post-result decoder audit finds that activation RBF
exceeds E5-interaction RBF only in the coherence-matched within-control analysis
and does not beat specified text baselines in cross-control transfer.
Separability transfers more reliably than orientation; the Qwen3 ladder shows
no consistent trend across both controls.
Learned-direction steering specifically moves an independent readout, but two
of five orthogonal control directions produce larger self-report shifts, and no
predicted safety-refusal effect appears.

The evidence therefore reaches different levels: it is strong for
prompt-category decodability, limited for construct-level generalization,
positive for causal control of the independent readout, insufficient for
direction-specific causal control of evaluation self-report, and inconclusive
for task-policy use. Probe accuracy alone therefore does not establish
evaluation awareness or identify which abstraction the model represents. Probe-based auditing
remains promising, but stronger semantic claims require
controls that separate evaluation status from textual form, competitive text
baselines, generalization across independent constructions, and interventions
tied to behavior. A probe's claimed meaning should be no stronger than the
controls that identify it. More broadly, trustworthy operational decisions
should depend on what a diagnostic has been shown to measure, not merely on its
predictive accuracy.

\bibliographystyle{plainnat}
\bibliography{references}

\appendix

\section{Additional experimental details}
\label{app:details}

\subsection{Dataset construction and sampling}

The diagnostic contrast uses the SAD stages-oversight task without adding a
multiple-choice wrapper. Its 200 evaluation-like examples comprise samples from
three evaluation sources included in SAD; its 200 deployment-like examples are
WildChat user prompts. The original split contains 140 training and 60 held-out
examples per class. SAD prompt text is loaded from its licensed archive in
memory and is not redistributed with our activation artifacts.

The single-turn transfer grid contains six binary constructions. Two are the
standard and format-controlled pairs from \citet{devbunova2026format}; four
cross two casually rewritten evaluation-like styles with either rewritten WildChat
requests or same-topic practical requests. Each new cell contains 250 examples.
The initial source mixture is 95 ARC, 58 MMLU, and 97 GSM8K items. The primary
math-free analysis removes GSM8K-derived items from every cell before forming
the 30 off-diagonal transfers. This leaves 153 items (95 ARC and 58 MMLU) in
each newly written cell. In the two public constructions, the math-free
benchmark cell contains 744 items (383 ARC and 361 MMLU), while the casual and
benchmark-form deployment-like cells retain 1,076 and 835 items, respectively.

The retained-history control begins with 250 filtered WildChat conversations.
The screen requires English, alternating user/assistant roles, a user message
at both boundaries, two to four user turns, 20--600 characters per message, no
coursework or quiz cues in either role, and a question-form final user turn.
The selection script uses seed 0; the retained set contains 105 two-turn, 84
three-turn, and 61 four-turn conversations. Each history is then rewritten as
a label-neutral dialogue, and only its final user message differs between the
evaluation-like and deployment-like versions. Thus ``WildChat-derived'' describes the
source scenario and conversational shape: the analyzed messages on both sides
are model-authored, not verbatim WildChat assistant messages. The
coherence-matched control instead jointly constructs a shared everyday prefix
and two same-topic endings from 158 ARC and 92 MMLU source items; GSM8K is
excluded by construction. Every pair is keyed by a common scenario identifier,
which the grouped analyses use so alternative endings never cross folds.

Constructed cells were written in iterative, label-balanced interactive Claude
sessions. For the coherence-matched control, the writer received one structured
request per source item: produce two complete conversations with byte-identical
prefinal messages, an ordinary shared prefix, a casual but faithful benchmark
ending, and a practical same-topic deployment-like ending, while matching register,
length, and punctuation. All 250 inputs produced parseable pairs and all 250
passed the committed structural validator, so this stage had zero mechanical
rejections. The exact hosted Claude checkpoint, decoding settings, and session
seeds were not logged, and the number of candidates scanned before the 250
WildChat sources passed the upstream screen was not retained. Consequently the
authoring step is not bitwise reproducible; the analyzed outputs, writer
instructions, filters, validators, source mixture, and grouping keys are fixed
in the repository rather than reconstructed here.

Constructed single-turn records expose a stable \texttt{uid}, batch, and style;
paired multi-turn records expose \texttt{item\_id} (for example,
\texttt{sk\_0000}). These are the grouping keys used by every split and
bootstrap. The repository
(\url{https://github.com/groverneev/eval-awareness-probes}) releases all
constructed prompts and these
identifiers. Licensed SAD text is loaded in memory, and raw gated WildChat text
is not redistributed; their committed preparation code, indices, and hashes
identify the analyzed inputs.

The coherence-matched pairs additionally underwent an independent blinded
language-model audit. An isolated \texttt{gpt-5.6-luna} reviewer saw randomized
A/B endings without source identifiers or condition labels and rated coherence,
topic match, and naturalness. Both endings in all 250 pairs were rated coherent,
topic-matched, and at least 3/5 in naturalness; mean ratings were 4.28 for
evaluation-like endings and 4.22 for deployment-like endings. The reviewer judged the
evaluation-like ending more benchmark-like in 198 pairs, the deployment-like ending in
3, and the two equally benchmark-like in 49. Thus, the audit supports the
intended coherence and topic matching, while confirming that task-form cues
remain.

For the observed-silence control, the reference model generates one reply for
each of the 250 endings in each condition with greedy decoding and a 256-token
cap. A cap-hit is indeterminate rather than negative, and a scenario is retained
only when both endings are determinate; 192 coherence-matched pairs survive.
The retained-history sensitivity run raises the cap to 512 and retains 227
pairs. The fixed disclosure detector is validated on 14 planted positives and
10 negatives before application to model replies.

\subsection{Checkpoint and extraction details}

The reference checkpoint is Meta-Llama-3.1-8B-Instruct. Its early activation
arrays contain activations but not embedded provenance; the sole local cached
snapshot used by the project resolves to
\nolinkurl{0e9e39f249a16976918f6564b8830bc894c89659}. We therefore report that
revision as reconstructed from the run environment, not as a revision stored
inside the arrays. Later replication artifacts freeze Qwen/Qwen3-8B at
\nolinkurl{b968826d9c46dd6066d109eabc6255188de91218},
google/gemma-3-12b-it at
\nolinkurl{96b6f1eccf38110c56df3a15bffe176da04bfd80}, and
mistralai/Ministral-3-8B-Instruct-2512-BF16 at
\nolinkurl{f6fae9795746f63c9be8344932f01275f3c63734}. Dataset hashes, manifest
hashes, and adapter identities are embedded in those result artifacts.

Extraction uses unquantized bfloat16, batch size one, and each tokenizer's
native chat template with no added system message and
\texttt{add\_generation\_prompt=False}. No Qwen-specific
\texttt{enable\_thinking} override is passed; the pinned Qwen3 template's
native rendering is used, including its empty thinking block when applicable.
Gemma and Ministral use validated text-only adapters to their official
multimodal processors. Stored activations are converted to float32. The
prespecified primary readout is the residual-stream state at the token closing
the user turn. For layer selection, the embedding output is indexed as layer
zero.

The scale-ladder repository labels and exact parameter counts are Qwen/Qwen3-
0.6B (0.752B parameters), 1.7B (2.032B), 4B (4.022B), 8B (8.191B), 14B
(14.768B), and 32B (32.762B). Their exact revisions are:
\begin{center}
\footnotesize
\begin{tabular}{ll}
0.6B & \texttt{c1899de289a04d12100db370d81485cdf75e47ca} \\
1.7B & \texttt{70d244cc86ccca08cf5af4e1e306ecf908b1ad5e} \\
4B & \texttt{1cfa9a7208912126459214e8b04321603b3df60c} \\
8B & \texttt{b968826d9c46dd6066d109eabc6255188de91218} \\
14B & \texttt{40c069824f4251a91eefaf281ebe4c544efd3e18} \\
32B & \texttt{9216db5781bf21249d130ec9da846c4624c16137}
\end{tabular}
\end{center}
They are original dense Qwen3 checkpoints run on one A100-SXM4-80GB (driver
580.126.16) under Python
3.12.3, PyTorch 2.9.1+cu128, CUDA 12.8, and Transformers 5.15.0. Gemma and
Ministral replication runs used separate NVIDIA A40 48GB GPUs; the Qwen
replication ran locally. Sparse MoE checkpoints, later Qwen generations, and
quantized weights are excluded before execution. Exact software and GPU
metadata were not embedded in the early reference-model arrays.

\subsection{Compute resources}

Compute accounting combines retained pod records with reconstructed estimates
for early runs whose exact wall-clock logs were not preserved. The initial
reference-model extraction, diagnostic contrasts, paired controls, pilots,
failed controls, and reruns used a local RTX 5070 Ti 16GB, free 2$\times$T4
16GB workers, and rented GPUs with at least 24GB; together these early and
preliminary experiments used approximately 15 GPU-hours. The cross-family
panel used the local worker for Qwen and separate A40 48GB RunPod pods for Gemma and
Ministral. Their measured full extractions took 11m09s and 12m15s, followed by
approximately 4.5 and 5.3 minutes of CPU analysis, respectively; including
setup, downloads, smoke tests, and artifact transfer, the two A40 pods were
allocated for about two hours in total. The six-checkpoint Qwen3 ladder used one
A100-SXM4-80GB pod for approximately six wall-clock hours, including model
downloads, extraction, analysis, and artifact transfer. Generation-heavy
silence and steering runs used rented 24GB-or-larger GPUs, including an RTX
4090 24GB for the 840-generation orthogonal-control audit; these runs required
approximately 2--4 GPU-hours in aggregate. Cross-model pods provisioned 100GB
of disk, while large activation artifacts were copied to local storage after
validation.

Across all reported, pilot, failed, and preliminary work, we estimate roughly
40 allocated GPU-hours. CPU-only probe fitting, text baselines, bootstrap and
permutation inference, figure generation, and validation required an additional
approximately 10--15 CPU-hours on a local desktop. These figures are
project-level resource estimates rather than exact utilization measurements:
pod allocations include download and transfer overhead, and the total is not
the sum of fully saturated accelerator time.

\subsection{Asset licenses and terms}

Table~\ref{tab:asset-licenses} records the licenses or provider terms for the
third-party datasets, model weights, embedding model, and hosted model services
used in the study. Dataset and model creators are credited in the main text and
references. We use these assets subject to their upstream terms, preserve
required attribution and notices, and do not redistribute third-party model
weights, licensed SAD prompt text, or raw WildChat conversations.

\begin{table}[h]
  \caption{Third-party asset licenses and use restrictions. Exact model
  revisions are listed in Appendix~\ref{app:details}.}
  \label{tab:asset-licenses}
  \centering
  \scriptsize
  \setlength{\tabcolsep}{3.5pt}
  \begin{tabularx}{\linewidth}{p{0.21\linewidth}p{0.24\linewidth}Y}
    \toprule
    Asset & License or terms & Use and redistribution treatment \\
    \midrule
    SAD stages-oversight
      & \href{https://github.com/LRudL/sad/blob/main/LICENSE.txt}{CC BY 4.0}; upstream anti-contamination protocol
      & Loaded from the licensed archive in memory; prompt text is not redistributed. \\
    WildChat-1M
      & \href{https://huggingface.co/datasets/allenai/WildChat-1M/blob/main/LICENSE.md}{ODC Attribution 1.0} for the database; separate content and privacy rights may apply
      & Used under gated-access and privacy controls; raw conversations are not redistributed. \\
    ARC
      & \href{https://huggingface.co/datasets/allenai/ai2_arc}{CC BY-SA 4.0}
      & Source items are credited; constructed derivatives retain provenance. \\
    MMLU and GSM8K
      & MIT licenses (\href{https://github.com/hendrycks/test/blob/master/LICENSE}{MMLU}; \href{https://github.com/openai/grade-school-math/blob/master/LICENSE}{GSM8K})
      & Source items are credited; constructed derivatives retain provenance. \\
    Evaluation-awareness $2\!\times\!2$
      & \href{https://huggingface.co/datasets/viliana-dev/eval-awareness-2x2}{MIT}
      & Public comparison cells are used with attribution. \\
    Llama-3.1-8B-Instruct
      & \href{https://huggingface.co/meta-llama/Llama-3.1-8B-Instruct}{Llama 3.1 Community License}
      & Accessed under the model license; weights are not redistributed. \\
    Qwen3 checkpoints
      & \href{https://huggingface.co/Qwen/Qwen3-8B}{Apache 2.0}
      & Accessed under the model license; weights are not redistributed. \\
    Gemma-3-12B-IT
      & \href{https://huggingface.co/google/gemma-3-12b-it}{Gemma Terms of Use}
      & Accessed under the model terms; weights are not redistributed. \\
    Ministral-3-8B-Instruct
      & \href{https://huggingface.co/mistralai/Ministral-3-8B-Instruct-2512-BF16}{Apache 2.0}
      & Accessed under the model license; weights are not redistributed. \\
    E5-base-v2
      & \href{https://huggingface.co/intfloat/e5-base-v2}{MIT}
      & Frozen embeddings are computed locally; weights are not redistributed. \\
    Hosted Claude and OpenAI models
      & Applicable provider service terms
      & Used only for constructed-text authoring and blinded review; no hosted-model weights are accessed or redistributed. \\
    \bottomrule
  \end{tabularx}
\end{table}

For the replication suite, each user-turn boundary is rendered and processed
as an independent causal prefix. This was frozen after a full-sequence smoke
test exposed suffix-length-dependent bfloat16 differences at otherwise
identical earlier tokens. Rejected smoke caches are isolated from the full-run
artifacts. Each completed model processes all 250 scenarios under both paired
controls, and every pre-final equality check is performed before probe fitting.
The Llama reference caches predate this change and retain small pre-final
differences (Appendix~\ref{app:results}).

\subsection{Probe fitting and inference details}
\label{app:fitting-details}

Repeated-holdout estimates use five paired 70/30 splits. Direct probe--text and
cross-model comparisons use grouped five-fold out-of-fold predictions, with
both endings of each scenario in the same fold and grouped layer selection
nested inside each outer training fold. In the single-turn grid, each source
split fixes the probe layer and orientation before any target is scored. We
report signed target AUROC and the fraction of ordered transfers preserving
orientation in a majority or all five split seeds. Separation strength
$|\mathrm{AUROC}-0.5|$ is retained only as a secondary measure of target
separability regardless of orientation. A
standardized L2-logistic probe and full-pipeline label permutations test
estimator and layer-selection robustness.

The post-result decoder audit fits standardized L2-logistic and RBF-kernel
classifiers at each outer fold's primary-selected layer. Logistic
$C\in\{0.001,0.01,0.1,1\}$ and RBF scaling, $C\in\{0.1,1,10\}$, and a
dimension-normalized kernel-width grid are selected in the same pair-grouped
inner folds; neither learned estimator searches layers. The full grouped
analysis is repeated over fold seeds 0--4. For nonlinear transfer, layer,
orientation, and RBF hyperparameters are selected on the source control only
before fitting all source pairs and scoring the other control.

Within an inner fold, an exact score tie is resolved by the fixed declaration
order: the lower-index layer first, then the earlier hyperparameter tuple. The
strongest sparse-text baseline is likewise the highest-AUROC out-of-fold member
of the prespecified order (unigram counts, word TF--IDF, character TF--IDF,
format, then length), with the earlier member winning an exact tie. Selection
uncertainty is not propagated through the reported bootstrap intervals.

Text baselines include length, 13 structural features, unigram counts, word
1--2-gram TF--IDF, and character-within-word 3--5-gram TF--IDF. Frozen E5 views
include the ending, full visible conversation, and concatenated prefix and
ending embeddings with their absolute difference and elementwise product. The
capacity-matched comparison applies RBF classifiers to sparse TF--IDF and E5
views, including an exact comparison with the strongest E5 prefix--ending
interaction. Pair-cluster bootstrap intervals use 2,000 draws; paired
approximate randomization swaps prediction vectors within scenarios. The scale
audit uses 5,000 paired draws that resample common scenario IDs across sizes and
refit the OLS slope against $\log_{10}$ exact parameter count.

E5 uses \texttt{intfloat/e5-base-v2} at revision
\nolinkurl{f52bf8ec8c7124536f0efb74aca902b2995e5bcd}, the
\texttt{query:} prefix, explicit user/assistant role markers, mean pooling and
L2 normalization, and left truncation to 512 wordpieces. Its representations
and logistic readout ($C=1$, liblinear) were frozen before execution. The
prefix--ending interaction was promoted to the main full-input comparator after
the initial results because it was the strongest fixed E5 view; it was not
refit or retuned. No analysis was externally preregistered. Frozen repository
prespecifications distinguish confirmatory analyses from post-result audits.
Prespecified analyses are the primary paired-control probes, the single-turn
transfer grid, the three cross-model replications, the Qwen3 scaling criterion,
and the steering protocol. Post-result analyses are the full-input E5 audit, the
decoder-sensitivity (L2-logistic and RBF) audit, the steering robustness audit
(prompt-level inference and orthogonal directions), and the steering--readout
cosine.
The direct probe-over-baseline randomization test uses a one-sided alternative;
other reported intervals and tests are per endpoint. We report all tested
cells, directions, and steering outcomes without family-wise multiplicity
correction.

\subsection{Intervention and generation details}
\label{app:intervention-details}

Steering uses the reference Llama checkpoint and the coherence-matched
last-user activations. The direction is the unit difference between the
evaluation-like and deployment-like class means. At residual block output $L$, the hook
applies
$h_L\leftarrow h_L+\alpha\lVert\mu_{\mathrm{eval}}-
\mu_{\mathrm{deploy}}\rVert_2\hat d$; it acts once during the prompt pass and
once on each newly generated token under the key/value cache. Hidden-state
index $L$ corresponds to transformer block $L-1$, and $\alpha=0$ installs no
hook. Before behavioral generation, layer 14 was selected from layers
3, 7, 10, and 14 using an independent readout's monotonicity, Spearman
correlation, and endpoint span.

The frozen coefficient grid is
$\{-4,-2,-1,-0.5,0,0.5,1,2,4\}$. Generation is greedy
(\texttt{do\_sample=False}) with at most 256 new tokens. The 36 self-report
carriers append the identical question about whether the interaction is an
evaluation; the 48 safety carriers are fixed borderline requests. This yields
756 learned-direction generations. Self-report labels follow the first matched
evaluation- or real-interaction term within 400 characters after negation
handling; refusals require a decline phrase within the first 120 characters. An
endpoint fails quality if its repeated four-gram fraction exceeds 0.35 or its
mean output length is less than half the zero-steering mean; both $\pm4$
endpoints fail and are excluded from trend summaries.

Prompt-level inference uses the paired $+2$ minus $-2$ risk difference, a
10,000-draw prompt-cluster percentile bootstrap with seed 20260823, and a
two-sided exact McNemar/binomial test. Mean within-prompt slopes use the seven
nondegenerate coefficients and the same bootstrap. Five controls are generated
from NumPy PCG64 seeds 20260823--20260827, projected exactly orthogonal to the
learned unit direction, renormalized, and multiplied by the same class-mean-gap
scale. Each is evaluated at $\alpha\in\{-2,0,2\}$, with the common zero point
reused, for 840 additional generations. Control generation ran on an RTX 4090
24GB GPU; independent-readout scoring used CPU offload. Prompts, coefficients,
seeds, scoring rules, and per-prompt outputs are retained in the intervention
artifacts.

\section{Complete results}
\label{app:results}

\subsection{Leakage check and estimator robustness}

Table~\ref{tab:appendix-estimators} reports the second-estimator check. Both
linear estimators reproduce the final-turn effects. Pre-final rows report the
expected exact-equality check, not awareness evidence. The permutation test
reruns layer selection after each label shuffle; with 500 draws, its smallest
attainable value is $1/501$.

\begin{table}[h]
  \caption{Repeated-holdout AUROC (mean $\mathbin{\pm}$ standard deviation over
  five paired 70/30 splits). Pre-final rows are leakage checks; permutation
  $p$ corresponds to the direction-of-means pipeline.}
  \label{tab:appendix-estimators}
  \centering
  \small
  \begin{tabular}{lccc}
    \toprule
    Contrast/readout & Direction of means & L2 logistic & Permutation $p$ \\
    \midrule
    Diagnostic SAD/WildChat, final & $0.9918\pm0.0085$ & $0.9933\pm0.0077$ & 0.002 \\
    Retained-history, final & $0.9917\pm0.0033$ & $0.9976\pm0.0017$ & 0.002 \\
    Retained-history, pre-final & $0.4985\pm0.0011$ & $0.5001\pm0.0003$ & 0.990 \\
    Coherence-matched, final & $0.9513\pm0.0115$ & $0.9894\pm0.0081$ & 0.002 \\
    Coherence-matched, pre-final & $0.4990\pm0.0015$ & $0.4992\pm0.0024$ & 0.940 \\
    \bottomrule
  \end{tabular}
\end{table}

\begin{figure}[h]
  \centering
  \includegraphics[width=0.55\linewidth]{figures/leakage_check.pdf}
  \caption{Reference-model paired-control leakage check. Token-identical
  prefixes yield pre-final states equal up to bfloat16 differences from
  full-sequence extraction; pre-final AUROCs are at chance. Labels become decodable only after
  the differing final turn (mean $\mathbin{\pm}$ SD over five splits).}
  \label{fig:leakage-check}
\end{figure}

The full retained-history trajectory is $0.9917$ at the final turn and $0.4975$,
$0.5003$, and $0.5003$ at one, two, and three turns before the end, with 250,
250, 145, and 61 available pairs. Embedding-layer states are identical in every
pair. At later layers, 34 of 456 retained-history pre-final states (24 pairs)
and 14 of 250 coherence-matched pairs differ, by at most 10\% and 12\% in
relative norm. These differences arise from suffix-length-dependent bfloat16
arithmetic in full-sequence extraction, the effect that led us to extract the
replication suite with independent causal-prefix passes, and they explain why
pre-final AUROCs deviate slightly from exactly $0.5$.

\subsection{Cross-model paired-control comparisons}
\label{app:crossmodel}

\begin{table}[h]
  \caption{Grouped five-fold final-turn AUROC for the primary linear probe and
  full-input E5 prefix--ending interaction. $\Delta$ is probe minus E5; bold
  differences have 95\% pair-bootstrap intervals below zero.}
  \label{tab:appendix-crossmodel}
  \centering
  \small
  \begin{tabular}{lrrr@{\qquad}rrr}
    \toprule
      & \multicolumn{3}{c}{Retained-history} & \multicolumn{3}{c}{Coherence-matched} \\
    \cmidrule(lr){2-4}\cmidrule(lr){5-7}
    Model & Probe & Text & $\Delta$ & Probe & Text & $\Delta$ \\
    \midrule
    Llama-3.1-8B & 0.9878 & 0.9955 & $\mathbf{-0.0077}$ & 0.9468 & 0.9532 & $-0.0064$ \\
    Qwen3-8B & 0.9919 & 0.9955 & $-0.0036$ & 0.9450 & 0.9532 & $-0.0082$ \\
    Gemma-3-12B & 0.9730 & 0.9955 & $\mathbf{-0.0226}$ & 0.8590 & 0.9532 & $\mathbf{-0.0942}$ \\
    Ministral-3-8B & 0.9831 & 0.9955 & $\mathbf{-0.0124}$ & 0.9279 & 0.9532 & $\mathbf{-0.0253}$ \\
    \bottomrule
  \end{tabular}
\end{table}

\begin{table}[h]
  \caption{Activation-minus-specified-text AUROC margins by decoder capacity.
  Linear $\Delta$ uses the full-input E5 interaction with fixed L2 logistic
  regression. RBF $\Delta$ uses an RBF over that E5 view and reports the mean
  over five grouped-fold seeds. The columns therefore use different estimators
  over the same text representation.}
  \label{tab:appendix-nonlinear-margins}
  \centering
  \small
  \setlength{\tabcolsep}{7pt}
  \begin{tabular}{lrr@{\qquad}rr}
    \toprule
      & \multicolumn{2}{c}{Retained-history} &
        \multicolumn{2}{c}{Coherence-matched} \\
    \cmidrule(lr){2-3}\cmidrule(lr){4-5}
    Model & Linear $\Delta$ & RBF $\Delta$ & Linear $\Delta$ & RBF $\Delta$ \\
    \midrule
    Llama-3.1-8B & $-0.0077$ & $+0.0020$ & $-0.0064$ & $+0.0388$ \\
    Qwen3-8B & $-0.0036$ & $+0.0035$ & $-0.0082$ & $+0.0390$ \\
    Gemma-3-12B & $-0.0226$ & $-0.0071$ & $-0.0942$ & $+0.0348$ \\
    Ministral-3-8B & $-0.0124$ & $+0.0024$ & $-0.0253$ & $+0.0401$ \\
    \bottomrule
  \end{tabular}
\end{table}

\FloatBarrier
\subsection{Within-family Qwen3 scale ladder}
\label{app:qwen-scale}

\begin{table}[h]
  \caption{Qwen3 scale ladder. Linear entries are grouped-OOF probe minus the
  full-input E5 interaction; nonlinear entries are activation RBF minus
  E5-interaction RBF, mean over five grouped-fold seeds.}
  \label{tab:qwen-scale}
  \centering
  \scriptsize
  \setlength{\tabcolsep}{3.4pt}
  \begin{tabular}{lrrrr}
    \toprule
      & \multicolumn{2}{c}{Primary linear $\Delta$} &
        \multicolumn{2}{c}{Nonlinear mean $\Delta$} \\
    \cmidrule(lr){2-3}\cmidrule(lr){4-5}
    Exact size & Shared & Coherent & Shared & Coherent \\
    \midrule
    0.75B & $-0.0283$ & $-0.0120$ & $-0.0024$ & $+0.0408$ \\
    2.03B & $-0.0072$ & $-0.0146$ & $+0.0024$ & $+0.0406$ \\
    4.02B & $-0.0085$ & $-0.0131$ & $+0.0017$ & $+0.0386$ \\
    8.19B & $-0.0035$ & $-0.0066$ & $+0.0035$ & $+0.0390$ \\
    14.77B & $-0.0429$ & $+0.0007$ & $+0.0036$ & $+0.0411$ \\
    32.76B & $+0.0009$ & $-0.0066$ & $+0.0027$ & $+0.0386$ \\
    \bottomrule
  \end{tabular}
\end{table}

The paired-bootstrap slopes are $+0.0054$
$[-0.0017,+0.0129]$ (retained-history) and $+0.0068$
$[+0.0013,+0.0132]$ (coherence-matched). Subtracting the same E5 vector at each
size leaves the prespecified slopes unchanged; they fail the frozen criterion
because only one interval excludes zero. For secondary sparse-text transfer,
shared-to-coherent has
slope $-0.0529$ $[-0.0664,-0.0389]$, whereas coherent-to-shared has slope
$-0.0002$ $[-0.0156,+0.0149]$. These six checkpoints test an internal-family
pattern for this pipeline, not a universal scaling law.

\begin{figure}[h]
  \centering
  \includegraphics[width=\linewidth]{figures/qwen_scale_ladder.pdf}
  \caption{Within-family Qwen3 scale ladder. \textbf{(a)} Primary linear probe
  minus the full-input E5 interaction, with 95\% pair-bootstrap intervals and
  fitted slopes. Using one common text vector changes margin levels but not the
  prespecified slopes; only the coherence-matched interval excludes zero.
  \textbf{(b)} Activation RBF minus E5-interaction RBF, mean $\mathbin{\pm}$ SD
  over five grouped-fold seeds, remains nearly flat with size.}
  \label{fig:qwen-scale}
\end{figure}

\FloatBarrier
\subsection{Cross-model readout sensitivity}
\label{app:readout-sensitivity}

\begin{table}[h]
  \caption{Uniform readout audit using the same paired folds and nested grouped
  layer selection. Entries are probe AUROCs. ``Token'' is the prespecified state
  at the token closing the final user turn. ``Mean (all)'' averages the final
  user turn and permits embedding layer zero; ``Mean ($L\geq1$)'' excludes that
  layer. The unchanged strongest-text AUROCs are 0.9529 (retained-history) and
  0.9447 (coherence-matched).}
  \label{tab:appendix-readout}
  \centering
  \scriptsize
  \setlength{\tabcolsep}{2.8pt}
  \begin{tabular}{lrr@{\qquad}rrr}
    \toprule
      & \multicolumn{2}{c}{Retained-history} & \multicolumn{3}{c}{Coherence-matched} \\
    \cmidrule(lr){2-3}\cmidrule(lr){4-6}
    Model & Token & Mean & Token & Mean (all) & Mean ($L\geq1$) \\
    \midrule
    Llama-3.1-8B & 0.9878 & 0.9916 & 0.9468 & 0.9322 & 0.9325 \\
    Qwen3-8B & 0.9919 & 0.9785 & 0.9450 & 0.9439 & 0.9332 \\
    Gemma-3-12B & 0.9730 & 0.9488 & 0.8590 & 0.9388 & 0.6380 \\
    Ministral-3-8B & 0.9831 & 0.9869 & 0.9279 & 0.8775 & 0.9380 \\
    \bottomrule
  \end{tabular}
\end{table}

Mean pooling is not uniformly better. Under coherence matching, permitting every
layer raises Gemma from 0.8590 to 0.9388 but lowers Ministral from 0.9279 to
0.8775; the direct mean-minus-token changes are $+0.0797$
($[+0.0563,+0.1037]$) and $-0.0504$ ($[-0.0710,-0.0288]$), respectively.
Crucially, Gemma and Qwen select layer zero in all five unrestricted-mean outer
folds (Llama in three; Ministral in one). Layer zero here is the mean of input
token embeddings, not a transformer-computed state. When it is excluded,
Llama, Qwen, and Ministral cluster at 0.9325--0.9380, whereas Gemma falls to
0.6380 and selects layers 4, 4, 48, 7, and 48; this unstable selection
indicates that no transformer layer yields a strong mean-pooled Gemma readout. The audit therefore does not
support a simple extraction-failure explanation for the primary Gemma result:
the apparent alignment under unrestricted mean pooling is carried by a lexical
embedding average, and the broader finding is architecture-by-readout
sensitivity. We retain the prespecified closing-token readout rather than
selecting a post-result readout for visual alignment.

\subsection{Same-layer decoder sensitivity}
\label{app:nonlinear-probe}

\begin{table}[h]
  \caption{Post-result same-layer decoder audit at the prespecified closing-token
  readout (fold seed 0). L2 and RBF hold each outer fold's primary-selected
  layer fixed while tuning their own hyperparameters in the same pair-grouped
  inner folds. Primary direction-of-means values appear in
  Table~\ref{tab:appendix-crossmodel}. $\Delta$ is activation RBF minus a
  capacity-matched RBF over the strongest frozen E5 prefix--ending interaction
  (0.9953 shared; 0.9566 coherent).}
  \label{tab:appendix-nonlinear}
  \centering
  \scriptsize
  \setlength{\tabcolsep}{2.5pt}
  \begin{tabular}{lrrr@{\qquad}rrr}
    \toprule
      & \multicolumn{3}{c}{Retained-history} & \multicolumn{3}{c}{Coherence-matched} \\
    \cmidrule(lr){2-4}\cmidrule(lr){5-7}
    Model & L2 & RBF & $\Delta$ E5 & L2 & RBF & $\Delta$ E5 \\
    \midrule
    Llama-3.1-8B & 0.9948 & 0.9978 & $+0.0025$ & 0.9872 & 0.9870 & $\mathbf{+0.0304}$ \\
    Qwen3-8B & 0.9983 & 0.9992 & $\mathbf{+0.0039}$ & 0.9863 & 0.9879 & $\mathbf{+0.0313}$ \\
    Gemma-3-12B & 0.9876 & 0.9899 & $-0.0055$ & 0.9757 & 0.9811 & $\mathbf{+0.0245}$ \\
    Ministral-3-8B & 0.9977 & 0.9988 & $\mathbf{+0.0034}$ & 0.9849 & 0.9871 & $\mathbf{+0.0305}$ \\
    \bottomrule
  \end{tabular}
\end{table}

The learned-linear gain is largest where the primary result was weakest. At
seed 0, coherence-matched L2-minus-direction-of-means is $+0.0404$ for Llama,
$+0.0412$ for Qwen, $+0.1166$ for Gemma, and $+0.0570$ for Ministral; all four
paired bootstrap intervals exclude zero. RBF-minus-L2 is instead $-0.0002$
$[-0.0034,+0.0027]$, $+0.0016$ $[-0.0009,+0.0044]$, $+0.0055$
$[-0.0004,+0.0123]$, and $+0.0022$ $[-0.0011,+0.0053]$, respectively. Thus
most of the earlier RBF-minus-direction-of-means gain is not specifically
nonlinear. Against the stronger E5 RBF, the coherence-matched margins and intervals
are respectively $+0.0304$ $[+0.0149,+0.0461]$, $+0.0313$
$[+0.0164,+0.0465]$, $+0.0245$ $[+0.0100,+0.0388]$, and $+0.0305$
$[+0.0158,+0.0459]$. On the retained-history control, only Qwen and Ministral have positive
intervals; Llama's interval crosses zero and Gemma is numerically below E5 with
an interval that also crosses zero.

Repeating the complete nested analysis over fold seeds 0--4 leaves the
estimator pattern intact. Mean coherence-matched L2 AUROC is 0.9875, 0.9864,
0.9740, and 0.9864 for Llama, Qwen, Gemma, and Ministral; mean RBF-minus-L2 is
$+0.0009$, $+0.0023$, $+0.0104$, and $+0.0033$. Mean activation RBF AUROC
ranges from 0.9844--0.9897 (SD 0.0015--0.0023), and every model exceeds the
fixed E5 RBF in every seed; mean margins are $+0.0388$, $+0.0390$, $+0.0348$,
and $+0.0401$. Retained-history margins remain stable but
heterogeneous: $+0.0020$, $+0.0035$, $-0.0071$, and $+0.0024$.

\begin{table}[h]
  \caption{Nonlinear cross-control transfer, mean $\mathbin{\pm}$ SD over five
  source-selection fold seeds. Activation RBF layer, orientation, and kernel
  settings are selected on the source only. The nested nonlinear text row
  likewise selects sparse TF--IDF or a frozen E5 view on the source only. The
  final row gives the stronger existing linear character-TF--IDF comparator.}
  \label{tab:appendix-nonlinear-transfer}
  \centering
  \small
  \begin{tabular}{lrr}
    \toprule
    Method/model & Shared $\rightarrow$ coherent & Coherent $\rightarrow$ shared \\
    \midrule
    Llama RBF & $0.772\mathbin{\pm}0.016$ & $0.708\mathbin{\pm}0.005$ \\
    Qwen RBF & $0.806\mathbin{\pm}0.014$ & $0.819\mathbin{\pm}0.013$ \\
    Gemma RBF & $0.715\mathbin{\pm}0.011$ & $0.809\mathbin{\pm}0.000$ \\
    Ministral RBF & $0.717\mathbin{\pm}0.000$ & $0.842\mathbin{\pm}0.006$ \\
    Nonlinear text & $0.827\mathbin{\pm}0.000$ & $0.830\mathbin{\pm}0.030$ \\
    Linear character TF--IDF & 0.866 & 0.849 \\
    \bottomrule
  \end{tabular}
\end{table}

Nonlinear activation transfer is below the source-selected nonlinear text
pipeline for every model in the shared-to-coherent direction and does not
reliably exceed it in the reverse direction; linear character TF--IDF is
stronger still. The near-ceiling within-control gains therefore do not reveal a
construction-invariant decision rule. Because the audit was motivated by the
observed direction-of-means results, holds that estimator's chosen layer fixed,
and is largely reproduced by L2 logistic, it supports estimator-sensitive
decodability rather than specifically nonlinear geometry, a post-hoc
replacement of the primary estimator, or a semantic identification claim.

\subsection{Reference-model sparse-text comparisons and transfer}

\begin{figure}[h]
  \centering
  \includegraphics[width=\linewidth]{figures/disjoint_transfer.pdf}
  \caption{Reference-model (Llama-3.1-8B-Instruct) cross-construction transfer
  against unigram counts on the math-free grid. The source fixes which score
  direction means evaluation-like.
  \textbf{(a)} Seed-averaged signed target AUROC for every ordered probe
  transfer; values below $0.5$ reverse orientation, and hatched cells share a
  constituent dataset cell. \textbf{(b)} Mean signed AUROC
  ($\mathbin{\pm}$ split-seed SD) and the number of transfers preserving
  orientation in a majority of seeds. Absolute distance from chance is retained
  as a secondary separability diagnostic in the text.}
  \label{fig:transfer}
\end{figure}

The secondary grouped sparse-text audit gives probe-minus-TF--IDF differences
$+0.0349$ $[+0.0183,+0.0540]$ on retained history and $+0.0021$
$[-0.0134,+0.0168]$ under coherence matching. Shared-to-coherent and
coherent-to-shared transfer differences are $-0.0212$
$[-0.0489,+0.0081]$ and $-0.1298$ $[-0.1665,-0.0908]$, respectively.

The observed-silence analysis retains 192 coherence-matched pairs; the probe
and unigram AUROCs are $0.9369$ and $0.9354$. The retained-history 512-token
sensitivity run retains 227 pairs and yields $0.9908$ versus $0.9460$. In every
tested cell, the detector finds zero positive awareness disclosures among 250
generated replies; planted-control validation catches all 14 positives with no
false alarms among ten negatives.

For steering, the selected layer is 14. Across nondegenerate coefficients
$-2,-1,-0.5,0,0.5,1,2$, evaluation self-report rates are respectively
$0.139,0.611,0.667,0.778,0.806,0.889,0.889$. Safety-refusal rates over the same
grid are $0.083,0.104,0.104,0.125,0.104,0.083,0.062$. Coefficients $\pm4$ fail
the frozen repetition/degeneracy rule and are excluded from both trends.

\begin{figure}[h]
  \centering
  \includegraphics[width=\linewidth]{figures/steering_effects.pdf}
  \caption{Reference-model (Llama-3.1-8B-Instruct) learned-direction steering
  curves. \textbf{(a)} Independent readout.
  \textbf{(b)} Evaluation self-report rate. \textbf{(c)} Safety-refusal rate.
  Shaded, hollow endpoints violate the frozen output-degeneracy rule and are
  excluded from trends. Orthogonal controls establish readout specificity but
  not self-report specificity (text and Table~\ref{tab:steering-controls}).}
  \label{fig:steering}
\end{figure}

\subsection{Reference-model prompt-level steering inference and orthogonal controls}
\label{app:steering-controls}

A post-result robustness audit treats prompt index as the sampling unit.
For the learned direction, the paired $+2$ minus $-2$ self-report risk difference
is $+0.750$ ($[+0.583,+0.889]$; exact McNemar $p=1.12\times10^{-7}$), while the
safety difference is $-0.021$ ($[-0.063,0.000]$; $p=1.0$). Mean within-prompt
slopes over all seven nondegenerate coefficients are $+0.176$
($[+0.136,+0.213]$) and $-0.006$ ($[-0.017,+0.002]$), respectively.

Each control is a seeded standard-normal layer-14 vector projected exactly
orthogonal to the learned unit direction, renormalized, and multiplied by the
same class-mean-gap scale. All ten control endpoints pass the frozen quality
screen. Table~\ref{tab:steering-controls} reports the paired effects.

\begin{table}[h]
  \caption{Paired $+2$ minus $-2$ risk differences for the learned and
  norm-matched orthogonal directions. Intervals are 95\% prompt bootstraps.}
  \label{tab:steering-controls}
  \centering
  \small
  \begin{tabular}{lrr}
    \toprule
    Direction & Self-report difference & Safety difference \\
    \midrule
    Learned & $+0.750$ $[+0.583,+0.889]$ & $-0.021$ $[-0.063,0.000]$ \\
    Seed 20260823 & $+0.417$ $[+0.250,+0.583]$ & $-0.042$ $[-0.104,0.000]$ \\
    Seed 20260824 & $+0.944$ $[+0.861,+1.000]$ & $-0.021$ $[-0.104,+0.063]$ \\
    Seed 20260825 & $+0.306$ $[+0.167,+0.472]$ & $0.000$ $[-0.083,+0.083]$ \\
    Seed 20260826 & $0.000$ $[-0.083,+0.083]$ & $-0.042$ $[-0.104,0.000]$ \\
    Seed 20260827 & $+1.000$ $[+1.000,+1.000]$ & $-0.042$ $[-0.104,0.000]$ \\
    \bottomrule
  \end{tabular}
\end{table}

The learned direction moves the independent readout by $+1.045$ class-gap
units, versus $-0.094$ to $+0.084$ for the controls. It therefore passes
readout specificity but fails the frozen behavioral-specificity criterion:
two orthogonal directions cause larger self-report shifts. The result supports
causal accessibility of the decoded readout, not a direction-specific causal
link to self-report.

\section{Failed and rejected controls}
\label{app:failures}

The experimental sequence contained several probe-blind failures. Reporting
them matters because each failure removes an interpretation that a high AUROC
would otherwise invite.

\paragraph{Single-turn surface gate.}
Length matching succeeds in the four new cells, but format and lexical matching
do not. Five-fold out-of-fold surface AUROCs are:

\begin{table}[h]
  \caption{Surface gate for the four new single-turn contrasts. EL denotes the
  two evaluation-like styles, WC the WildChat-derived deployment-like cell, and
  TM the topic-matched deployment-like cell.}
  \label{tab:appendix-surface-gate}
  \centering
  \small
  \begin{tabular}{lrrrr}
    \toprule
    Contrast & Character length & Word length & Format & Unigram \\
    \midrule
    EL A vs. WC & 0.4726 & 0.4932 & 0.6627 & 0.9319 \\
    EL A vs. TM & 0.4764 & 0.4673 & 0.8326 & 0.9783 \\
    EL B vs. WC & 0.4923 & 0.5222 & 0.7818 & 0.9056 \\
    EL B vs. TM & 0.4964 & 0.4954 & 0.7435 & 0.9135 \\
    \bottomrule
  \end{tabular}
\end{table}

\paragraph{Leaking multi-turn construction.}
The first written ``disguised evaluation'' cell obtains final-boundary probe
AUROC $0.9846$, but its probe is already at $0.9492$ after the first user turn,
before the benchmark question appears. Its synthetic histories therefore leak
the label and fail the pre-final check. We retain the result only
as the failure that motivated the retained-history recombination.
Figure~\ref{fig:leak} contrasts the two constructions.

\begin{figure}[ht]
  \centering
  \includegraphics[width=0.60\linewidth]{figures/failed_temporal_control.pdf}
  \caption{The rejected leaking construction and repaired control.
  \textbf{(a)} Written histories are authored per label, whereas retained-history
  pairs are byte-identical before the final turn. \textbf{(b)} Reference-model
  probe AUROC (mean $\mathbin{\pm}$ standard deviation over five splits). Both
  rejected cells are near ceiling pre-final and fail the check; retained history
  passes the leakage check.}
  \label{fig:leak}
\end{figure}

\paragraph{Rejected coherence attempt.}
An initial attempt appended unrelated benchmark items to deployment-like histories
but failed manual coherence review. We report only the jointly constructed
replacement.

\paragraph{Authorship-factorial pilot.}
A prespecified pilot independently varied evaluation status and history author.
All 150 histories terminated, but the frozen blinded model and human gate
retained 36 of 50 scenarios, below the 40-scenario activation threshold. We
extracted or inspected no activations and make no history-authorship invariance
claim. Because the gate was frozen before any activation was examined, running
the pilot on the 36 retained scenarios would be a post-result analysis; we leave
an adequately powered authorship factorial to future work.

\clearpage
\end{document}
